\documentclass[11pt]{article}

\usepackage[margin=1in]{geometry}
\usepackage{amsmath,amssymb,amsthm}
\usepackage{hyperref}

\newtheorem{theorem}{Theorem}[section]
\newtheorem{lemma}[theorem]{Lemma}
\newtheorem{proposition}[theorem]{Proposition}
\newtheorem{corollary}[theorem]{Corollary}
\newtheorem{definition}[theorem]{Definition}
\newtheorem{conjecture}[theorem]{Conjecture}
\newtheorem{remark}[theorem]{Remark}

\newcommand{\R}{\mathbb{R}}
\newcommand{\one}{\mathbf{1}}
\newcommand{\cut}{\operatorname{cut}}
\newcommand{\vol}{\operatorname{vol}}
\newcommand{\diag}{\operatorname{diag}}
\newcommand{\Sc}{\operatorname{Sc}}
\newcommand{\E}{\mathcal{E}}
\newcommand{\Null}{\operatorname{null}}
\newcommand{\polylog}{\operatorname{polylog}}

\title{Spectral Maximizers and Cut-Preserving Rewiring}
\author{Working manuscript}
\date{\today}

\begin{document}
\maketitle

\begin{abstract}
Spectral clustering is limited by the possible gap between conductance and
the second generalized eigenvalue.  The 2019 spectral modification framework
proposed to close this gap by replacing an input graph by a cut-similar graph
whose low spectrum is much larger.  This note gives a theorem-driven version
of that program.  For a connected weighted graph $G$ and a terminal-Steiner
hierarchical cut approximator $T$, we define the spectral maximizer
$H_T$ as the Schur complement of $T$ onto the original terminal set.
We prove that $H_T$ is cut-similar to $G$, spectrally dominates every graph
whose cuts are comparable to those of $G$ up to the unavoidable diameter
factor of the hierarchy, and satisfies a generalized Cheeger inequality
linear in the quality and terminal diameter of $T$.  We then separate the
algorithmic story into three implementable layers: exact implicit application
of $H_T$ through tree solves, sparse edge-addition versions via Schur
complement sparsification, and practical tree-pack variants that generalize
the existing energy-based modification heuristic.  The final section connects
this construction to graph rewiring for graph neural networks, especially
spectral-gap and effective-resistance rewiring.
\end{abstract}

\section{Introduction}

Let $G=(V,E_G,w_G)$ be a connected weighted undirected graph.  Spectral
partitioning solves a relaxation whose value is the second eigenvalue of a
normalized or generalized Laplacian.  The target combinatorial objective is a
cut ratio.  Cheeger's inequality guarantees that these two quantities are
related, but only quadratically in the worst case.  This gap is not an artifact:
families of graphs are known where the eigenvector-driven separator is
asymptotically worse than the optimal cut \cite{guattery1998quality}.

The spectral modification proposal of \cite{koutis2019spectral} is to leave
the cut structure essentially intact while replacing the graph by a graph with
larger low eigenvalues.  The conceptual construction is a hierarchical
cut approximator $T$ with original vertices as terminals and additional
Steiner vertices.  Eliminating the Steiner vertices by a Schur complement
produces a dense terminal graph $H_T$.  Equivalently, one can keep the
Steiner representation and apply $H_T$ implicitly by solving a Dirichlet
problem on $T$.

This manuscript makes the following points precise.

\begin{itemize}
\item A terminal-Steiner hierarchy $T$ induces a maximizer
$H_T=\Sc(T,V)$.  The graph $H_T$ has both a Steiner-tree representation and
an edge-addition representation through star-mesh elimination.
\item If $T$ is an $\alpha$-cut approximator with terminal diameter $h$, then
$H_T$ is $\alpha h$-cut similar to $G$.  Moreover, for every terminal graph
$F$ whose cuts are bounded by $\rho$ times the cuts of $G$, the domination
$F \preceq h\rho H_T$ holds.
\item For every connected demand graph $B$, the pair $(H_T,B)$ satisfies
\[
  \lambda_2(H_T,B) \le \phi(H_T,B)
  \le \alpha h \lambda_2(H_T,B).
\]
The proof explicitly carries the path-support factor $h$.
\item Known hierarchical congestion approximator constructions of Racke,
Peng, and later work give existential and algorithmic maximizers with
polylogarithmic factors \cite{racke2002minimizing,racke2008optimal,
racke2014computing,peng2016approximate,henzinger2025improved,
li2025bottomup,agarwal2024parallel}.
\item Sparse Schur-complement approximations convert the dense maximizer into
a sampled edge set of near-linear size while preserving the guarantees up to
$(1\pm\varepsilon)$ factors \cite{lee2015sparsified,li2018spectral}.
\end{itemize}

\section{Preliminaries}

\subsection{Graphs, cuts, and domination}

All graphs in this note are finite, undirected, and weighted with nonnegative
edge weights.  For a graph $X=(U,E_X,w_X)$, its Laplacian is $L_X$ and its
Dirichlet energy is
\[
  \E_X(x)=x^T L_X x
  =\sum_{\{u,v\}\in E_X} w_X(u,v)(x_u-x_v)^2 .
\]
For $S\subseteq U$, let $\mathbf{1}_S$ be its indicator and
\[
  \cut_X(S)=\E_X(\mathbf{1}_S).
\]
For two graphs $X$ and $Y$ on the same vertex set we write
$X\preceq Y$ if $L_X\preceq L_Y$, equivalently
$\E_X(x)\le \E_Y(x)$ for all $x$.

\begin{definition}[One-sided and two-sided cut comparability]
Let $X$ and $Y$ be graphs on the same vertex set.  We say that $X$ is
$\rho$-cut dominated by $Y$ if
\[
  \cut_X(S)\le \rho\,\cut_Y(S)\qquad\text{for all }S\subseteq V.
\]
We say that $X$ and $Y$ are $\rho$-cut similar if each graph is
$\rho$-cut dominated by the other.
\end{definition}

The one-sided notion is the one needed for spectral maximality: if every cut
of $F$ is no larger than $\rho$ times the corresponding cut of $G$, then the
maximizer for $G$ spectrally dominates $F$ up to the hierarchy diameter.

\subsection{Generalized conductance and generalized eigenvalues}

Let $A$ and $B$ be connected graphs on the same vertex set $V$.  The graph
$B$ is the demand graph; it determines which notion of volume or balance is
being used.  Define
\[
  \phi(A,B)=\min_{\emptyset\ne S\subsetneq V}
    \frac{\cut_A(S)}{\cut_B(S)} .
\]
Since $B$ is connected, $\cut_B(S)>0$ for every nontrivial cut.

The second generalized eigenvalue of the pair $(A,B)$ is
\[
  \lambda_2(A,B)=
  \min_{x\notin \operatorname{span}\{\one\}}
  \frac{\E_A(x)}{\E_B(x)} .
\]
This quotient is invariant under adding a constant to $x$.  Equivalently,
one may minimize over any fixed complement of $\operatorname{span}\{\one\}$.
The first generalized eigenvalue is zero, with eigenvector $\one$.

\begin{lemma}[Cut upper bounds eigenvalue]
\label{lem:lambda-le-phi}
For connected $A$ and $B$,
\[
  \lambda_2(A,B)\le \phi(A,B).
\]
\end{lemma}

\begin{proof}
For every nontrivial $S$, the vector $\mathbf{1}_S$ is not constant modulo
$\operatorname{span}\{\one\}$ and
\[
  \frac{\E_A(\mathbf{1}_S)}{\E_B(\mathbf{1}_S)}
  =\frac{\cut_A(S)}{\cut_B(S)} .
\]
Taking the minimum over $S$ gives the claim.
\end{proof}

\begin{definition}[Demand complete graph]
Let $d\in\R_{>0}^V$ and $D=\sum_v d_v$.  The demand complete graph $K_d$ has
edge weight $d_ud_v/D$ on each unordered pair $\{u,v\}$.  Then
\[
  L_{K_d}=\diag(d)-dd^T/D
\]
and
\[
  \cut_{K_d}(S)=\frac{d(S)(D-d(S))}{D}.
\]
\end{definition}

If $d$ is the degree vector of $A$, then $\lambda_2(A,K_d)$ is the usual
second eigenvalue of the normalized Laplacian of $A$.  The conductance
$\phi(A,K_d)$ is the product-normalized cut objective.  It is within a factor
of two of the common $\cut_A(S)/\min\{\vol_A(S),\vol_A(\bar S)\}$ convention.

\section{Terminal-Steiner Hierarchies}

\begin{definition}[Terminal-Steiner hierarchy]
\label{def:hierarchy}
Let $G=(V,E_G,w_G)$ be connected.  A terminal-Steiner hierarchy for $G$ is a
weighted rooted tree
\[
  T=(V\dot\cup I,E_T,w_T,r)
\]
with terminal set $V$ and internal or Steiner set $I$, together with a set
$S_a\subseteq V$ for every node $a\in V\dot\cup I$, satisfying:
\begin{enumerate}
\item $S_r=V$.
\item The terminal leaves are identified with singletons: $S_v=\{v\}$ for
each $v\in V$.
\item If $a$ has children $b_1,\ldots,b_k$, then
$S_{b_1},\ldots,S_{b_k}$ are disjoint and their union is $S_a$.
\item If $a\ne r$ and $p(a)$ is the parent of $a$, then
\[
  w_T(a,p(a))=\cut_G(S_a).
\]
\end{enumerate}
The terminal diameter $h(T)$ is the maximum number of tree edges on a path
between two terminals.
\end{definition}

The canonical weights in item 4 imply an important path-lifting identity.  For
each edge $e=\{u,v\}\in E_G$, let $P_e$ be the unweighted terminal path
between $u$ and $v$ in $T$, and let $w_eP_e$ denote the tree graph on this
path with every path edge given weight $w_G(e)$.

\begin{lemma}[Path lifting identity]
\label{lem:path-lifting}
For every terminal-Steiner hierarchy $T$ of $G$,
\[
  T=\sum_{e\in E_G} w_G(e)P_e
\]
as weighted graphs on $V\dot\cup I$.
\end{lemma}

\begin{proof}
Fix a tree edge $(a,p(a))$.  The path $P_{\{u,v\}}$ contains this tree edge
if and only if exactly one of $u,v$ lies in $S_a$.  Therefore the total
weight contributed to $(a,p(a))$ by the right-hand side is precisely
$\cut_G(S_a)$, which equals $w_T(a,p(a))$ by Definition
\ref{def:hierarchy}.
\end{proof}

\begin{definition}[Cut approximation quality]
\label{def:cut-approx}
For $S\subseteq V$, define the best binary extension cut
\[
  \mu_T(S)=\min_{Y\subseteq I}\cut_T(S\cup Y).
\]
The hierarchy $T$ is an $\alpha$-cut approximator for $G$ if for every
$S\subseteq V$,
\[
  \cut_G(S)\le \mu_T(S)\le \alpha\,\cut_G(S).
\]
\end{definition}

The lower bound follows automatically from the path-lifting identity for the
canonical hierarchies above, but it is useful to include it as part of the
interface.  The upper bound is the nontrivial approximation guarantee supplied
by hierarchical cut/congestion approximator algorithms.

\section{The Spectral Maximizer}

Order the vertices of $T$ so that internal vertices $I$ come first and
terminals $V$ come last.  Write
\[
  L_T=
  \begin{pmatrix}
    L_{II} & L_{IV}\\
    L_{VI} & L_{VV}
  \end{pmatrix}.
\]
Since $T$ is connected and every component of $T[I]$ touches a terminal
through the boundary, the Dirichlet block $L_{II}$ is positive definite.

\begin{definition}[Spectral maximizer]
\label{def:maximizer}
The spectral maximizer induced by $T$ is the terminal graph $H_T$ whose
Laplacian is the Schur complement
\[
  L_{H_T}=L_{VV}-L_{VI}L_{II}^{-1}L_{IV}.
\]
\end{definition}

\begin{lemma}[The maximizer is connected]
\label{lem:max-connected}
The matrix $L_{H_T}$ is the Laplacian of a connected weighted graph on $V$.
Equivalently, $\Null(L_{H_T})=\operatorname{span}\{\one\}$.
\end{lemma}

\begin{proof}
Schur complements of Laplacians onto a nonempty terminal set are Laplacians
with nonpositive off-diagonal entries and zero row sums.  It remains to check
the nullspace.  If $\E_{H_T}(x)=0$, Lemma \ref{lem:harmonic} gives an
extension $z=(y,x)$ with $\E_T(z)=0$.  Since $T$ is connected, every edge of
$T$ has equal endpoint values under $z$, so $z$ is constant and therefore
$x$ is constant.  The reverse inclusion is immediate from zero row sums.
\end{proof}

\begin{lemma}[Harmonic extension]
\label{lem:harmonic}
For every terminal vector $x\in\R^V$,
\[
  \E_{H_T}(x)
  =
  \min_{y\in\R^I}
  \E_T(y,x).
\]
The minimizer is the unique internal vector
$y_x=-L_{II}^{-1}L_{IV}x$.
\end{lemma}

\begin{proof}
The energy of $T$ at fixed terminal values $x$ is the quadratic
\[
  y^T L_{II}y+2y^TL_{IV}x+x^TL_{VV}x .
\]
Its unique minimizer satisfies $L_{II}y+L_{IV}x=0$.  Substituting the
minimizer gives the Schur complement quadratic form.
\end{proof}

\begin{proposition}[Two representations]
\label{prop:representations}
The maximizer $H_T$ may be used in either of the following equivalent ways.
\begin{enumerate}
\item Steiner representation: store $T$ and apply $L_{H_T}$ by the formula
\[
  L_{H_T}x=L_{VV}x-L_{VI}L_{II}^{-1}L_{IV}x .
\]
\item Edge-addition representation: eliminate vertices of $I$ one at a time.
Eliminating a Steiner vertex $q$ with incident conductances
$c_1,\ldots,c_k$ to neighbors $a_1,\ldots,a_k$ and
$c_q=\sum_i c_i$ deletes $q$ and adds a clique on its neighbors with edge
conductance $c_ic_j/c_q$ between $a_i$ and $a_j$.  After all Steiner
vertices are eliminated, the resulting terminal Laplacian is $L_{H_T}$.
\end{enumerate}
\end{proposition}

\begin{proof}
The first item is the block formula for the Schur complement.  The second is
the standard star-mesh formula for eliminating one Laplacian variable.  Schur
complements are associative, so eliminating the vertices of $I$ in any order
gives the same terminal Schur complement.
\end{proof}

\section{Spectral Maximality}

The following elementary path-support inequality is the source of the
terminal-diameter loss.

\begin{lemma}[Path support]
\label{lem:path-support}
Let $P$ be a path of length $\ell$ between endpoints $s,t$, with every path
edge having conductance $w$.  Let $E_{st}$ be the single edge of conductance
$w$ between $s$ and $t$.  Then
\[
  E_{st}\preceq \ell P .
\]
\end{lemma}

\begin{proof}
For any vector $z$ on the path,
\[
  w(z_s-z_t)^2
  =
  w\left(\sum_{i=1}^{\ell}(z_{i-1}-z_i)\right)^2
  \le
  \ell w\sum_{i=1}^{\ell}(z_{i-1}-z_i)^2
\]
by Cauchy-Schwarz.
\end{proof}

\begin{lemma}[The hierarchy supports the input graph]
\label{lem:G-supported}
If $T$ is a terminal-Steiner hierarchy for $G$ and $h=h(T)$, then
\[
  G\preceq hT
\]
where $G$ is viewed as a graph on $V\dot\cup I$ with isolated Steiner
vertices.
\end{lemma}

\begin{proof}
By Lemma \ref{lem:path-lifting}, $T=\sum_{e\in E_G} w_eP_e$.  For each
$e=\{u,v\}$, the corresponding terminal path has length at most $h$, so
$w_eE_{uv}\preceq h\,w_eP_e$ by Lemma \ref{lem:path-support}.  Summing over
$e$ gives the claim.
\end{proof}

\begin{theorem}[Conditional spectral maximizer theorem]
\label{thm:conditional-max}
Let $T$ be an $\alpha$-cut approximating terminal-Steiner hierarchy for a
connected weighted graph $G$, let $h=h(T)$, and let $H_T=\Sc(T,V)$.
Then:
\begin{enumerate}
\item $H_T$ and $G$ are $\alpha h$-cut similar in the explicit sense
\[
  \frac{1}{\alpha}\cut_{H_T}(S)\le \cut_G(S)
  \le h\,\cut_{H_T}(S)
  \qquad\text{for all }S\subseteq V.
\]
\item If $F$ is any graph on $V$ with
$\cut_F(S)\le \rho\,\cut_G(S)$ for all $S\subseteq V$, then
\[
  F\preceq h\rho\,H_T .
\]
\end{enumerate}
\end{theorem}

\begin{proof}
For the upper comparison $\cut_{H_T}(S)\le \alpha\cut_G(S)$, use Lemma
\ref{lem:harmonic} with $x=\mathbf{1}_S$:
\[
  \cut_{H_T}(S)
  =
  \min_{y\in\R^I}\E_T(y,\mathbf{1}_S)
  \le
  \min_{Y\subseteq I}\E_T(\mathbf{1}_Y,\mathbf{1}_S)
  =
  \mu_T(S)
  \le
  \alpha\cut_G(S).
\]

For the lower comparison $\cut_G(S)\le h\cut_{H_T}(S)$, Lemma
\ref{lem:G-supported} gives, for every internal extension $y$,
\[
  \cut_G(S)=\E_G(\mathbf{1}_S)
  \le h\,\E_T(y,\mathbf{1}_S).
\]
Taking the minimum over $y$ and applying Lemma \ref{lem:harmonic} gives the
claim.

Now let $F$ be $\rho$-cut dominated by $G$.  Build the path lift
\[
  T_F=\sum_{\{u,v\}\in E_F} w_F(u,v)P^T_{uv},
\]
where $P^T_{uv}$ is the terminal path between $u$ and $v$ in $T$.  For the
tree edge $(a,p(a))$, the weight in $T_F$ is $\cut_F(S_a)$, hence it is at
most $\rho\cut_G(S_a)=\rho w_T(a,p(a))$.  Therefore $T_F\preceq \rho T$.
By path support, $F\preceq hT_F\preceq h\rho T$ on $V\dot\cup I$.  Thus for
all terminal vectors $x$ and all internal extensions $y$,
\[
  \E_F(x)\le h\rho\,\E_T(y,x).
\]
Minimize over $y$ and use Lemma \ref{lem:harmonic} to obtain
$F\preceq h\rho H_T$.
\end{proof}

\begin{remark}
The word ``maximizer'' is not meant to assert uniqueness.  It says that,
after fixing a cut hierarchy $T$, the Schur complement $H_T$ spectrally
dominates every graph with comparable cuts, up to the terminal diameter and
cut-comparison factors.  Different hierarchies can produce different
maximizers.
\end{remark}

\section{Existence from Hierarchical Cut Approximators}

The previous section is conditional on the hierarchy.  Existing work on
oblivious routing and congestion approximators supplies such hierarchies.
The exact constants and logarithmic powers depend on the chosen black box.

\begin{theorem}[HCA black box]
\label{thm:hca-blackbox}
For every connected weighted graph $G$ on $n$ vertices, there is a
terminal-Steiner hierarchy $T$ with terminal diameter $h=O(\log n)$ and
cut-approximation quality $\alpha=\Gamma(n)$, where $\Gamma(n)$ is
polylogarithmic.  Moreover:
\begin{enumerate}
\item Peng's nearly-linear max-flow framework gives a near-linear-time
construction with polylogarithmic $\Gamma(n)$
\cite{peng2016approximate}.
\item Recent hierarchical congestion approximators improve the best
near-linear quality to $O(\log^2 n\log\log n)$ with high probability, and
also admit parallel implementations with near-linear work
\cite{henzinger2025improved,agarwal2024parallel}.
\item Bottom-up constructions give an alternative nearly-linear route to
polylogarithmic congestion approximators without the same recursive
alternation between max-flow and approximator construction
\cite{li2025bottomup}.
\end{enumerate}
\end{theorem}

This theorem is used as a black box here.  Its objects are often presented as
hierarchical congestion approximators or oblivious routing trees rather than
as the terminal-Steiner hierarchy in Definition \ref{def:hierarchy}.  The
translation is standard: the laminar family of cuts becomes the rooted tree,
and each child-parent edge is assigned the capacity of the corresponding cut
in $G$.

\begin{corollary}[Existence of spectral maximizers]
\label{cor:existence}
Every connected weighted graph $G$ has a spectral maximizer $H$ that is
$O(\Gamma(n)\log n)$-cut similar to $G$ and that spectrally dominates every
graph $F$ with $F$ $\rho$-cut dominated by $G$ up to an
$O(\rho\log n)$ factor.  With the current near-linear HCA bound quoted above,
the cut-similarity factor is
$O(\log^3 n\log\log n)$ with high probability.
\end{corollary}

\begin{proof}
Apply Theorem \ref{thm:hca-blackbox} and then Theorem
\ref{thm:conditional-max}.
\end{proof}

\section{Cheeger Inequalities for the Maximizer}

\begin{theorem}[Generalized Cheeger inequality for $H_T$]
\label{thm:cheeger}
Let $T$ be an $\alpha$-cut approximating terminal-Steiner hierarchy for $G$,
let $h=h(T)$, and let $H_T=\Sc(T,V)$.  For every connected demand graph
$B$ on $V$,
\[
  \lambda_2(H_T,B)\le \phi(H_T,B)
  \le \alpha h\,\lambda_2(H_T,B).
\]
\end{theorem}

\begin{proof}
The lower inequality $\lambda_2(H_T,B)\le \phi(H_T,B)$ is Lemma
\ref{lem:lambda-le-phi}.  We prove the other direction.

Let $x$ be a minimizer for $\lambda_2(H_T,B)$ in the quotient by constants,
and let $z=(y_x,x)$ be its harmonic extension to $T$ as in Lemma
\ref{lem:harmonic}.  Thus
\[
  \lambda_2(H_T,B)=\frac{\E_T(z)}{\E_B(x)} .
\]
For each edge $e=\{u,v\}\in E_B$, let $P^T_{uv}$ be the terminal path in
$T$.  By Lemma \ref{lem:path-support}, and because $z$ has terminal values
$x$,
\[
  w_B(e)(x_u-x_v)^2
  \le h\, w_B(e)\E_{P^T_{uv}}(z).
\]
Summing over $e\in E_B$ gives
\[
  \E_B(x)\le h\,\E_{T_B}(z),
  \qquad
  T_B=\sum_{e\in E_B} w_B(e)P^T_e .
\]
Therefore
\[
  \lambda_2(H_T,B)
  \ge
  \frac{1}{h}\frac{\E_T(z)}{\E_{T_B}(z)}
  \ge
  \frac{1}{h}
  \min_{a\ne r}
  \frac{w_T(a,p(a))}{w_{T_B}(a,p(a))}.
\]
The last inequality is the elementary fact that a quotient of weighted sums
is at least the minimum quotient of corresponding coefficients; terms with
zero denominator are ignored.

For a tree edge $(a,p(a))$, the denominator is
$w_{T_B}(a,p(a))=\cut_B(S_a)$, while the numerator is
$w_T(a,p(a))=\cut_G(S_a)$.  From Theorem
\ref{thm:conditional-max}, $\cut_{H_T}(S_a)\le \alpha\cut_G(S_a)$, hence
\[
  \frac{w_T(a,p(a))}{w_{T_B}(a,p(a))}
  =
  \frac{\cut_G(S_a)}{\cut_B(S_a)}
  \ge
  \frac{1}{\alpha}
  \frac{\cut_{H_T}(S_a)}{\cut_B(S_a)} .
\]
Taking the minimum over hierarchy clusters can only increase the minimum over
all nontrivial subsets of $V$.  Thus
\[
  \lambda_2(H_T,B)
  \ge
  \frac{1}{\alpha h}\phi(H_T,B),
\]
which is equivalent to the desired upper Cheeger bound.
\end{proof}

\begin{corollary}[Normalized cut specialization]
\label{cor:normalized}
Let $d\in\R_{>0}^V$ and let $K_d$ be the demand complete graph.  Then
\[
  \lambda_2(H_T,K_d)\le \phi(H_T,K_d)
  \le \alpha h\,\lambda_2(H_T,K_d).
\]
If $d$ is the degree vector of $H_T$, the eigenvalue is the standard second
normalized Laplacian eigenvalue of $H_T$.  If $d$ is the degree vector of
$G$, the same inequality holds for the original-volume normalization used by
the modification algorithm.
\end{corollary}

\begin{corollary}[Actual Cheeger inequality from spectral closeness]
\label{cor:actual}
Let $M$ be a terminal graph such that
\[
  c_- H_T\preceq M\preceq c_+ H_T
\]
for constants $0<c_-\le c_+$.  Then for every connected demand graph $B$,
\[
  \lambda_2(M,B)\le \phi(M,B)
  \le \frac{c_+}{c_-}\,\alpha h\,\lambda_2(M,B).
\]
\end{corollary}

\begin{proof}
The lower inequality is again Lemma \ref{lem:lambda-le-phi}.  For the upper
bound,
\[
  \phi(M,B)\le c_+\phi(H_T,B)
  \le c_+\alpha h\,\lambda_2(H_T,B)
  \le \frac{c_+}{c_-}\alpha h\,\lambda_2(M,B).
\]
\end{proof}

\section{Algorithms}

This section gives the algorithmic interface promised by the theory.  It
separates exact implicit maximizers from sparse sampled edge sets and from
practical tree-pack heuristics.

\subsection{Exact implicit maximizer}

\paragraph{Input.}
A connected weighted graph $G$ and a hierarchical cut approximator routine
returning $T=(V\dot\cup I,E_T,w_T,r)$.

\paragraph{Data structure.}
Store the four Laplacian blocks $L_{II},L_{IV},L_{VI},L_{VV}$ and a solver
for the tree Dirichlet block $L_{II}$.  Because $T$ is a tree, exact linear
time elimination is possible after rooting; in a numerical implementation,
a sparse Cholesky or specialized tree solve is enough.

\paragraph{Matrix-vector product.}
For $x\in\R^V$:
\begin{enumerate}
\item Compute $b=L_{IV}x$.
\item Solve $L_{II}y=b$.
\item Return $L_{VV}x-L_{VI}y$.
\end{enumerate}
This is exactly $L_{H_T}x$.  It supports generalized eigensolvers for
$L_{H_T}x=\lambda L_Bx$ without materializing dense terminal edges.

\paragraph{Public interface.}
\[
\begin{array}{ll}
\texttt{build\_maximizer(G)} & \text{return hierarchy, blocks, and solver},\\
\texttt{apply\_maximizer(state,x)} & \text{return }L_{H_T}x,\\
\texttt{eigs\_maximizer(state,B,k)} & \text{return the first }k\text{ generalized eigenvectors.}
\end{array}
\]

\subsection{Sparse edge sampler}

The dense edge-addition representation can be approximated using Schur
complement sparsification.  The needed theorem is now standard in the
Laplacian-solver literature \cite{lee2015sparsified,li2018spectral}.

\begin{theorem}[Sparse Schur-complement sampler, black-box form]
\label{thm:schur-sampler}
Given a Laplacian $L_T$ and terminal set $V$, there is a randomized algorithm
that returns a terminal graph $\widehat H$ with
\[
  (1-\varepsilon)H_T\preceq \widehat H\preceq(1+\varepsilon)H_T
\]
with high probability and with
$O(|V|\polylog |V|/\varepsilon^2)$ terminal edges, up to the precise
polylogarithmic factors of the chosen Schur-sparsifier routine.
\end{theorem}

\begin{corollary}[Sparse sampled maximizers inherit Cheeger]
\label{cor:sparse-cheeger}
If $\widehat H$ satisfies Theorem \ref{thm:schur-sampler} with
$0<\varepsilon<1$, then for every connected demand graph $B$,
\[
  \lambda_2(\widehat H,B)\le\phi(\widehat H,B)
  \le
  \frac{1+\varepsilon}{1-\varepsilon}\,
  \alpha h\,\lambda_2(\widehat H,B).
\]
Moreover $\widehat H$ is cut-similar to $G$ with the factors in Theorem
\ref{thm:conditional-max} multiplied by at most $1+\varepsilon$ or
$(1-\varepsilon)^{-1}$.
\end{corollary}

\begin{proof}
The proof is the same as Corollary \ref{cor:actual}, with
$c_-=1-\varepsilon$ and $c_+=1+\varepsilon$.  Spectral approximation implies
cut approximation by testing indicator vectors.
\end{proof}

\paragraph{Sampling by star-mesh elimination.}
An implementable route is to eliminate Steiner vertices one at a time.  When
eliminating a vertex of degree $k$, the exact star-mesh clique has
$\binom{k}{2}$ edges with weights $c_ic_j/c_q$.  Instead of inserting the
whole clique, sample edges from the clique with probabilities proportional to
leverage-score or effective-resistance surrogates, reweighting sampled edges
unbiasedly.  The analysis should follow the martingale matrix concentration
template used in Schur-complement sparsifiers.  For a first prototype, one
can use oversampling probabilities proportional to $c_ic_j/c_q$ and then
measure the actual spectral error empirically.

\paragraph{Public interface.}
\[
\begin{array}{ll}
\texttt{sample\_maximizer\_edges(G,budget,eps,seed)}
 & \text{return sampled weighted edges plus certificates},\\
\texttt{certificate.cut}
 & \text{sampled cut tests against }G\text{ and }H_T,\\
\texttt{certificate.spectral}
 & \text{Ritz/eigenvalue tests for }(1\pm\varepsilon)\text{ approximation}.
\end{array}
\]

\subsection{Adding sampled edges back to the original graph}

There are two natural uses of the sampled maximizer.

\paragraph{Replacement mode.}
Run spectral clustering on $\widehat H$ or on its implicit Steiner version.
The cuts returned by sweep or rounding are evaluated back in $G$.  Theorem
\ref{thm:conditional-max} and Corollary \ref{cor:sparse-cheeger} certify that
cuts in $\widehat H$ correspond to cuts in $G$ within polylogarithmic factors.

\paragraph{Augmentation mode.}
Use
\[
  M_\eta=G+\eta\widehat H
\]
for some $\eta>0$.  Then $M_\eta$ preserves the original graph channel while
adding long-range maximizer edges.  Since $G\preceq hH_T$ and
$\widehat H\preceq(1+\varepsilon)H_T$,
\[
  M_\eta\preceq (h+\eta(1+\varepsilon))H_T.
\]
Also $M_\eta\succeq \eta(1-\varepsilon)H_T$.  Therefore Corollary
\ref{cor:actual} gives the conservative certificate
\[
  \phi(M_\eta,B)
  \le
  \frac{h+\eta(1+\varepsilon)}{\eta(1-\varepsilon)}
  \alpha h\,\lambda_2(M_\eta,B).
\]
This bound is not meant to be optimized; it shows that augmentation restores a
linear Cheeger relationship whenever the maximizer component is not too small.

\subsection{Tree-pack maximizers}

The MATLAB code accompanying \cite{koutis2019spectral} computes several
high-energy spanning trees, replaces each tree by a tree maximizer, and sums
the resulting Steiner components.  The following proposition captures what is
already provable for this class of algorithms.

\begin{proposition}[Certified tree-pack modification]
\label{prop:tree-pack}
Let $T_1,\ldots,T_k$ be weighted subtrees of $G$, each on terminal set $V$.
Let $R_j$ be an $\alpha_j$-cut approximating hierarchy for $T_j$ with
terminal diameter $h_j$, and let $H_j=\Sc(R_j,V)$.  For weights
$\eta_j\ge0$, define
\[
  M=G+\sum_{j=1}^k \eta_j H_j .
\]
Then:
\begin{enumerate}
\item $M\succeq G$, so $\lambda_i(M,B)\ge \lambda_i(G,B)$ for every connected
demand graph $B$ and every generalized eigenvalue index $i$.
\item For every cut $S$,
\[
  \cut_M(S)
  \le
  \left(1+\sum_{j=1}^k\eta_j\alpha_j\right)\cut_G(S).
\]
\end{enumerate}
\end{proposition}

\begin{proof}
The first item follows from positive semidefinite addition and the
Courant-Fischer min-max characterization.  For the second, Theorem
\ref{thm:conditional-max} applied to each subtree gives
$\cut_{H_j}(S)\le \alpha_j\cut_{T_j}(S)$.  Since $T_j$ is a weighted subgraph
of $G$, $\cut_{T_j}(S)\le\cut_G(S)$.  Summing proves the bound.
\end{proof}

The proposition is deliberately modest: it certifies monotone spectral
increase and bounded cut distortion, but it does not say that a small tree
pack approximates the full maximizer.  That is the key algorithmic conjecture.

\begin{conjecture}[Energy-guided tree-pack convergence]
\label{conj:tree-pack}
There is a polynomial or polylogarithmic number of iterations of an
energy-guided tree-pack algorithm such that the sum of tree maximizers
$\sum_j\eta_jH_j$ is a constant or polylogarithmic spectral approximation to
the full maximizer $H_T$ on the low-dimensional subspace relevant to the first
$k$ generalized eigenvectors.
\end{conjecture}

A provable version may need to sample trees from a low-stretch or oblivious
routing distribution rather than selecting only maximum-energy spanning trees.
The existing energy method remains the right practical baseline: it directly
targets the Rayleigh quotient directions in which the current graph is far
from its maximizer.

\section{Connection to Graph Rewiring in GNNs}

Graph neural network rewiring tries to improve message passing by changing
the graph topology.  FoSR adds edges to increase spectral expansion and reduce
oversquashing, while using architectural separation to control oversmoothing
\cite{karhadkar2023fosr}.  Yusu Wang and coauthors analyze oversquashing
through effective resistance and add edges chosen to reduce total effective
resistance \cite{black2023oversquashing}.  A recent survey frames these
methods as part of a broad tradeoff between information flow and
oversmoothing \cite{attali2026survey}.

Spectral maximizers provide a complementary principle.  Instead of asking only
which edges most increase $\lambda_2$ or most reduce resistance, we ask for a
rewired graph whose cuts remain comparable to the original graph while its
low spectrum is raised as much as the cut structure permits.  This gives a
possible certificate missing from many rewiring heuristics: the added edges
improve long-range communication while preserving every cut up to a known
factor.

For GNNs, the clean experimental form is augmentation mode:
\[
  A_{\mathrm{message}} = A_G + \eta A_{\widehat H},
\]
with separate edge types or channels for original and maximizer edges.  The
separate-channel architecture is important: it lets the model distinguish
original geometry from cut-preserving shortcuts.  The mathematical certificate
is the cut/spectral approximation of $\widehat H$; the empirical question is
whether this reduces oversquashing without inducing oversmoothing.

\begin{conjecture}[Cut-certified rewiring for message passing]
For graph families where oversquashing is caused by high-diameter bottlenecks
rather than label-noise artifacts, sparse maximizer edges reduce effective
resistance and increase task-relevant spectral gaps while preserving the
class-relevant cut structure better than generic spectral-gap rewiring.
\end{conjecture}

\section{Numerical Checks and Benchmarks}

The first implementation should test the following cases.

\begin{enumerate}
\item Path graph.  The binary-tree hierarchy should raise $\lambda_2$ from
$\Theta(n^{-2})$ to near $\Theta((n\log n)^{-1})$, matching the example in
\cite{koutis2019spectral} up to logarithmic factors.
\item Expander graph.  The maximizer should not substantially alter the
useful low spectrum; this checks that the method does not damage graphs that
already have good Cheeger behavior.
\item Double-binary-tree and product bad cases.  These should reproduce the
known spectral separator failure and show improvement after modification
\cite{guattery1998quality,koutis2019spectral}.
\item Stochastic block and degree-corrected block models.  Compare replacement
and augmentation modes against regularized spectral clustering
\cite{zhang2018regularized}.
\item GNN rewiring tasks.  Compare against FoSR and effective-resistance
rewiring on graph classification benchmarks, reporting accuracy, resistance
statistics, spectral gap, and cut-distortion certificates.
\end{enumerate}

\section{Open Proof Tasks}

The core conditional theory above is complete once the HCA black box is
accepted.  The remaining proof work is algorithmic.

\begin{enumerate}
\item Instantiate Theorem \ref{thm:hca-blackbox} with a concrete HCA theorem
and write the exact quality/runtime statement, including probability of
success and hierarchy size.
\item Choose a Schur-complement sparsifier and replace Theorem
\ref{thm:schur-sampler} by its exact edge count and runtime.
\item Prove or refute Conjecture \ref{conj:tree-pack}.  A plausible route is
to sample trees from an oblivious-routing distribution and show that the
expected tree-pack maximizer spectrally approximates $H_T$ on a target
subspace.
\item Develop a posteriori certificates for practical edge samplers: sampled
cut tests, low-dimensional spectral approximation tests, and monotonicity of
generalized eigenvalues.
\end{enumerate}

\bibliographystyle{alpha}
\bibliography{references}

\end{document}
